% \section*{Broader Impact}
% The two techniques we developed in  this paper, i.e. the disentangled attention and the enhanced mask decoder, could be  extended to other related model structures, such as cross-attention, encoder-decoder attention and even the vanilla fully connected network. It has high potential to bring additional gains to a wide range of tasks built on top of these model structures.  Meanwhile, these two techniques could be applied to many fields beyond Natural Language Understanding (NLU), such as machine translation, spoken language understanding, text summarization and multi-modal pre-training with both text and image. On the other hand, the DeBERTa model proposed in this paper could be used to improve various Natural Language Processing(NLP) applications, such as intent identification in business documents or emails, conversation understanding in a custom support chatbot, query understanding in industrial search engines and fake news detection in social network.  

% DeBERTa is a pre-trained language model, similar to
% %GPT-2 \citep{radford2019gpt2},
% %GPT-3 \citep{brown2020language}, 
% BERT \citep{devlin2018bert}, 
% RoBERTa \citep{liu2019roberta}, 
% XLNet \citep{yang2019xlnet}, 
% UniLM \citep{dong2019unified},
% GPT \citep{radford2019language},
% ELMo \citep{peters2018elmo},
% and ELECTRA \citep{clark2020electra}. There are societal impacts introduced by a language model. 
% %T5 \citep{raffel2019t5}
% A better language model could bring a lot of positive societal impacts. For example, it can be used to better understand human language to build various business applications, such as a personal assistant like Amazon Alexa and Apple Siri. It can also improve human daily productivity to make life better. For instance, it can be used to build an application to perform automatic typo correction or synonym suggestion  while we are writing a document to make our text more professional.
% %that we can save a lot of time in typing text. 
% On the other hand, training a language model needs GPU resource and this will generate CO2 impacts. Although this  CO2 impact is applicable to most deep learning research today, we still need to be very careful and conservative of building very huge language models, such as the Google T5 (11B parameters)\citep{raffel2019t5} and the Open AI GPT-3(170B parameters)\citep{brown2020language}. If large companies such as Google, Facebook, Microsoft and Amazon can have a plan to achieve carbon neutral or negative, it would lead to huge positive environmental impacts.   

% Another under-explored impact related to most existing language models is the predictive bias when predicting a token. For example, for the sentence ``[MASK] is a doctor'', a language model may predict the [MASK] token to be `she' or `he' with different probabilities. This may create a gender bias. Generally, this bias is not specific to or introduced by a model structure itself. Instead, the model learned these probabilities from existing data. Any bias in the model could be a reflection of existing biases in our society and history. Therefore, it remains an open and important question on how to mitigate various biases in the future work, even though this is already out of the scope of this work. 


